% !TeX root = ../main.tex
\chapter{GIỚI THIỆU}\label{chap:introduction}

\section{Đặt vấn đề}

\subsection{Bối cảnh thực tiễn và bài toán cần giải quyết}
Trong hoạt động thương mại trực tuyến của các doanh nghiệp vừa và nhỏ kinh doanh sản phẩm vật lý thông qua quảng cáo và tư vấn, quá trình tạo giá trị là một chuỗi vận hành liên tục từ tiếp thị thu hút khách hàng, tiếp nhận và tư vấn hội thoại đa kênh, xác lập đơn hàng cho đến bàn giao thực thi nội bộ. Qua khảo sát trường hợp ban đầu do chủ dự án cung cấp tại một doanh nghiệp kinh doanh cây ăn quả và cây cảnh quy mô toàn quốc (gọi là Doanh nghiệp A), doanh nghiệp vừa phải triển khai nội dung truyền thông trên nhiều kênh trực tuyến, vừa tiếp nhận tin nhắn khách hàng, tư vấn bán hàng, ghi nhận đơn và chuyển giao cho bộ phận đóng gói, gửi hàng. Khi các công đoạn này vận hành trên những công cụ rời rạc hoặc trao đổi thủ công, doanh nghiệp phải đối mặt với bốn thách thức cốt lõi:

\begin{enumerate}
  \item \textbf{Sự đứt gãy giữa hoạt động tiếp thị đầu phễu và quy trình tư vấn -- đơn hàng:} Hoạt động sản xuất nội dung truyền thông, quản trị trang mạng xã hội và thu hút tương tác thường diễn ra trên các công cụ tách rời khỏi hệ thống bán hàng. Khi khách hàng phản hồi từ một chiến dịch tiếp thị, ngữ cảnh nguồn khách dễ bị thất lạc ngay từ điểm chạm đầu tiên; đồng thời, doanh nghiệp thiếu một nền tảng hạt nhân thống nhất về định danh tổ chức, phân quyền thành viên, quản lý kết nối kênh và truyền thông sự kiện để liên kết mở rộng sang các bài toán vận hành tiếp thị về sau.
  \item \textbf{Áp lực phản hồi lặp lại và rủi ro khi ứng dụng trí tuệ nhân tạo thiếu kiểm soát:} Khách hàng thường đặt các câu hỏi lặp lại về quy cách, giá bán, chính sách hoặc hỏi gộp nhiều nội dung. Ngay cả khi sử dụng lệnh gọi mẫu trả lời nhanh có sẵn, nhân viên tư vấn vẫn phải đọc ngữ cảnh, lựa chọn thủ công từng mẫu và chỉnh sửa lại câu chữ. Ngược lại, việc tích hợp mô hình ngôn ngữ lớn một cách tự phát mà không kiểm soát dữ liệu và trạng thái tiềm ẩn hai rủi ro lớn: mô hình tự suy đoán thông tin khi dữ liệu doanh nghiệp bị khuyết thiếu, hoặc gửi phản hồi tự động chồng chéo sau khi nhân viên đã chủ động tiếp quản hội thoại.
  \item \textbf{Khoảng trống trạng thái giữa bản nháp giao dịch và trách nhiệm bàn giao nội bộ:} Việc tách rời kênh nhắn tin và nơi ghi nhận đơn (thường bằng bảng tính rời rạc và báo miệng) khiến đơn hàng trao đổi dở dang không được lưu thành thực thể có trạng thái bền vững qua nhiều lượt tư vấn, thiếu cơ chế kiểm chứng điều kiện nghiệp vụ bắt buộc trước khi xác nhận đơn chính thức, và không phân định rõ thời điểm chuyển giao trách nhiệm giữa hình thức \emph{chuyển trực tiếp} và hình thức \emph{chờ xác nhận nhận} tùy theo tổ chức nhân sự của từng doanh nghiệp.
  \item \textbf{Yêu cầu cô lập dữ liệu và bảo toàn tính nhất quán trên nền tảng đa doanh nghiệp:} Người quản lý thiếu góc nhìn hợp nhất có thể đối chiếu ngược về bản ghi gốc để theo dõi tiến độ công việc, người chịu trách nhiệm hiện tại và nhu cầu sản phẩm. Trên một nền tảng phục vụ đồng thời nhiều doanh nghiệp, xác thực đăng nhập không đủ để bảo đảm an toàn dữ liệu; mọi thao tác đọc, ghi, xác nhận đơn, bàn giao công việc hay áp dụng đề xuất tự động đều phải tuân thủ phạm vi doanh nghiệp, vai trò, người phụ trách và phiên bản trạng thái hiện hành.
\end{enumerate}

\subsection{Hiện trạng các giải pháp tương tự trên thị trường và những đánh đổi hiện hữu}
Qua khảo sát tài liệu công khai và đặc tả kỹ thuật do nhà cung cấp công bố đối với các nền tảng quản trị và bán hàng phổ biến tại Việt Nam, thị trường hiện có ba nhóm giải pháp chính đã tích hợp dữ liệu và AI ở các mức độ khác nhau, song vẫn tồn tại những đánh đổi nhất định trước bài toán nêu trên:

\begin{itemize}
  \item \textbf{Nhóm nền tảng quản lý bán lẻ -- kho và trợ lý giao dịch (Sapo, KiotViet / KiCi, Knote):} Nhóm này giải quyết tốt nghiệp vụ danh mục sản phẩm, bán lẻ và tồn kho đa kênh. Trợ lý Sabi AI của Sapo tra cứu dữ liệu cửa hàng theo quyền của tài khoản đang đăng nhập, điều hướng màn hình và tạo bản nháp sản phẩm để chủ cửa hàng kiểm tra, xác nhận trước khi lưu~\cite{sapo-sabi,sapo-restaurant-ai}; KiCi hỗ trợ tra cứu số liệu kinh doanh bằng ngôn ngữ tự nhiên~\cite{kiotviet-kici}; Knote hỗ trợ hộ kinh doanh lên đơn bằng giọng nói, tra cứu doanh thu bằng ngôn ngữ tự nhiên và chụp hóa đơn chi phí~\cite{knote-vneconomy,knote-appstore}. Tuy nhiên, các bản nháp này được tạo theo từng đối tượng ngay tại màn hình quản trị và chờ người dùng xác nhận tức thời; nhiệm vụ báo cáo đặt lịch của KiCi hiện tính trên toàn bộ chi nhánh được phân quyền, chưa chọn riêng được một chi nhánh và bỏ qua kỳ không phát sinh dữ liệu~\cite{kiotviet-kici}. Quan trọng hơn, trọng tâm của nhóm này nằm ở chuẩn hóa nhập liệu tại điểm bán, chưa gắn kết trạng thái điều khiển lượt tư vấn với quy trình bàn giao trách nhiệm nhiều bước sau chốt đơn, cũng như chưa cung cấp kiến trúc sự kiện mở để kết nối với các ứng dụng vận hành tiếp thị chuyên biệt.
  \item \textbf{Nhóm nền tảng tập trung hội thoại đa kênh và kịch bản tự động (Pancake, Botcake):} Thế mạnh của nhóm này là gom hội thoại và bình luận từ nhiều kênh mạng xã hội, cung cấp lệnh trả lời nhanh, phân công cuộc hội thoại, kịch bản tự động và chatbot AI tự động trả lời câu hỏi thường gặp, đơn hàng và gợi ý sản phẩm~\cite{botcake-messenger}. Dẫu vậy, cơ chế phân công chỉ dừng ở mức cuộc hội thoại, chưa biểu diễn chuỗi bước công việc hậu cần với phiên bản bất biến và chế độ chờ xác nhận nhận. Bên cạnh đó, khảo sát đặc tả API và webhook công khai của Pancake~\cite{pancake-api} chưa ghi nhận giao diện lập trình mở để hệ thống nghiệp vụ ngoài lưu bản nháp có kiểm chứng vào vùng duyệt của nhân viên, đồng thời việc tích hợp cần mã truy cập riêng cho từng trang và chịu hạn mức 5 lệnh gọi/giây/trang.
  \item \textbf{Nhóm nền tảng quản trị doanh nghiệp tổng thể và quy trình động (MISA AMIS, 1Office, Base.vn):} Hệ sinh thái của nhóm này bao phủ rộng nhiều mảng nghiệp vụ: MISA AMIS là hệ sinh thái hơn 40 ứng dụng cùng trợ lý AI AVA~\cite{misa-amis,misa-ava}; 1Office cung cấp công cụ thiết kế không cần mã nguồn và tác nhân AI, với dịch vụ tư vấn triển khai được giới thiệu là phù hợp với doanh nghiệp trên 50 nhân sự~\cite{1office-ai}; Base.vn cho phép trợ lý AI ngoài kết nối qua giao thức ngữ cảnh mô hình (MCP) và hành động trong phạm vi quyền của tài khoản Base, kèm lớp chặn thao tác phá hủy ở phía máy chủ~\cite{base-mcp}. Đánh đổi đặt ra đối với doanh nghiệp thương mại quy mô nhỏ là độ phức tạp thiết lập ban đầu khi phải lựa chọn, liên kết nhiều phân hệ hành chính hoặc cấu hình kết nối trợ lý ngoài (ví dụ trình kết nối tùy chỉnh trên ChatGPT yêu cầu gói trả phí~\cite{base-mcp}), trong khi nhu cầu cấp thiết của nhóm doanh nghiệp này là một luồng vận hành gọn nhẹ xuyên suốt từ hội thoại tư vấn đến đơn nháp và bàn giao nội bộ.
\end{itemize}

Cần lưu ý rằng các nhận xét trên dựa trên nguồn tài liệu công khai, chưa thông qua đo đạc trực tiếp trên tài khoản thương mại của từng nhà cung cấp nên không dùng để khẳng định giới hạn tuyệt đối của toàn bộ hệ sinh thái đối thủ. Dẫu vậy, Bảng~\ref{tab:market-comparison} cho thấy một khoảng trống thiết kế rõ nét tại vùng giao thoa giữa mở rộng vận hành tiếp thị, hội thoại có kiểm soát AI, đơn nháp bền vững và quy trình bàn giao trách nhiệm.

\begin{table}[htbp]
  \centering
  \small
  \setlength{\tabcolsep}{4pt}
  \renewcommand{\arraystretch}{1.2}
  \caption{Đối chiếu các nhóm giải pháp hiện có và hướng tiếp cận của Hexta}\label{tab:market-comparison}
  \begin{tabularx}{\linewidth}{@{}
    >{\raggedright\arraybackslash}p{2.9cm}
    >{\raggedright\arraybackslash}X
    >{\raggedright\arraybackslash}X
    >{\raggedright\arraybackslash}X
    >{\raggedright\arraybackslash}X@{}}
    \toprule
    Nhóm giải pháp & Trọng tâm nghiệp vụ & Hội thoại và AI tư vấn & Đơn nháp và bàn giao & Mở rộng tiếp thị \\
    \midrule
    Bán lẻ -- kho và trợ lý giao dịch (Sapo, KiotViet/KiCi, Knote) &
    Quản lý bán lẻ, danh mục hàng hóa và tồn kho tại điểm bán. &
    Tra cứu tồn kho, hỏi đáp số liệu và hỗ trợ nhập liệu nhanh bằng ngôn ngữ tự nhiên. &
    Bản nháp AI tạo theo từng đối tượng, chờ xác nhận tại màn hình quản trị; chưa gắn quy trình bàn giao hậu cần đa bước. &
    Đóng khung trong tác nghiệp bán lẻ; chưa mở trục kết nối cho ứng dụng nội dung tiếp thị. \\
    Hội thoại đa kênh (Pancake, Botcake) &
    Gom tin nhắn, bình luận mạng xã hội và phản hồi trực tuyến. &
    Trả lời nhanh, kịch bản bot và chatbot AI; chưa mở giao diện lưu nháp cho hệ thống ngoài. &
    Phân công ở cấp cuộc hội thoại; chưa quản lý chuỗi bước công việc hậu cần có phiên bản. &
    Tập trung ở tầng tiếp nhận kênh; không quản lý vòng đời đơn hàng và điều phối nội bộ. \\
    Quản trị tổng thể (MISA AMIS, 1Office, Base.vn) &
    Hệ sinh thái đa phân hệ (nhân sự, tài chính, quy trình nội bộ). &
    Hỏi đáp báo cáo, tác nhân AI và trình kết nối trợ lý ngoài qua cấu hình phân quyền. &
    Quy trình tùy biến mạnh nhưng chi phí thiết lập cao đối với doanh nghiệp thương mại nhỏ. &
    Thiên về quản trị nội bộ; cần cấu hình phức tạp để gắn chặt trục tư vấn -- đơn hàng. \\
    \textbf{Định hướng của Hexta} &
    \textbf{Hạt nhân đa doanh nghiệp nối tư vấn, đơn hàng và điều phối việc.} &
    \textbf{Định tuyến kịch bản/AI chỉ dựa trên dữ liệu được cấp; khóa tự động khi người tiếp nhận.} &
    \textbf{Đơn nháp bền vững, cổng xác nhận chặt chẽ, quy trình bất biến hai chế độ bàn giao.} &
    \textbf{Thiết kế mở về định danh, quyền và sự kiện để kết nối các ứng dụng tiếp thị sau này.} \\
    \bottomrule
  \end{tabularx}
\end{table}

\subsection{Định hướng khác biệt và giá trị giải quyết của Hexta}
Khác với hướng tiếp cận xây dựng một hệ thống hoạch định nguồn lực đồ sộ hay một công cụ nhắn tin tách rời dữ liệu giao dịch, \textbf{Hexta} được thiết kế như nền tảng hạt nhân cho một hệ sinh thái đa ứng dụng hướng tới kết nối liền mạch giữa vận hành tiếp thị và vận hành bán hàng. Trong tầm nhìn mở rộng dài hạn, nền tảng lõi đóng vai trò trung tâm về định danh doanh nghiệp, phân quyền thành viên, tiếp nhận tương tác kênh, hỗ trợ tư vấn, quản lý đơn hàng và điều phối quy trình thực thi; từ đó tạo đường biên tích hợp sẵn sàng cho các ứng dụng chuyên biệt về khởi tạo nội dung truyền thông và quản trị xuất bản kênh tiếp thị trong các giai đoạn tiếp theo.

Trong phạm vi giải quyết trực tiếp khoảng trống giao thoa đã chỉ ra, Hexta tập trung vào năm giá trị thiết kế cốt lõi:
\begin{enumerate}
  \item \textbf{Liên kết liền mạch dữ liệu và định danh trên kiến trúc mở:} Kết nối xuyên suốt từ tiếp nhận hội thoại khách hàng (qua kênh trực tiếp hoặc môi trường mô phỏng có gắn nhãn), hồ sơ khách hàng, danh mục sản phẩm cho đến đơn hàng, công việc bàn giao và tổng hợp số liệu, đồng thời chuẩn hóa biên định danh doanh nghiệp và luồng sự kiện để sẵn sàng mở rộng hệ sinh thái đa ứng dụng.
  \item \textbf{Đơn nháp bền vững tách biệt với lệnh xác nhận có điều kiện:} Đơn nháp và đề xuất phản hồi được lưu trữ bền vững theo trạng thái và phiên bản qua nhiều lượt trao đổi; lệnh xác nhận đơn chỉ được chấp thuận khi máy chủ kiểm chứng đầy đủ thông tin khách hàng, người phụ trách, số lượng hợp lệ, đơn giá đã được cung cấp tường minh (tách biệt với giá bằng không) và địa chỉ theo nhu cầu giao hàng.
  \item \textbf{Quy trình công việc tuần tự có phiên bản bất biến và hai chế độ bàn giao:} Doanh nghiệp tùy chỉnh chuỗi bước xử lý theo mô hình tổ chức riêng và lựa chọn cơ chế chuyển trực tiếp hoặc chờ xác nhận nhận cho từng bước; công việc đã khởi tạo giữ cố định phiên bản quy trình ban đầu để tránh xáo trộn khi cấu hình thay đổi.
  \item \textbf{Ranh giới nghiêm ngặt giữa AI hỗ trợ và thẩm quyền nghiệp vụ:} Thành phần AI chỉ tiếp nhận ảnh chụp ngữ cảnh được cấp phép và trả về ứng viên phản hồi, không có quyền tự ghi dữ liệu nghiệp vụ hay phát lệnh gửi tin. Kết quả phải bám sát dữ liệu doanh nghiệp được cung cấp, tự động chuyển người xử lý khi thiếu thông tin, và bị vô hiệu hóa ngay lập tức theo hàng rào điều khiển khi nhân viên tiếp nhận hội thoại.
  \item \textbf{Nhất quán cạnh tranh, chống trùng lặp và truy vết đầy đủ:} Mọi thao tác thay đổi trạng thái đều kiểm tra phạm vi doanh nghiệp, đối chiếu phiên bản lạc quan để loại bỏ cập nhật dựa trên dữ liệu cũ, áp dụng cơ chế chống xử lý trùng lặp sự kiện và lưu vết lịch sử phục vụ đối chiếu minh bạch.
\end{enumerate}

\section{Mục tiêu và phạm vi nghiên cứu}
Mục tiêu tổng quát của khóa luận là thiết kế, xây dựng và đánh giá hệ thống hỗ trợ doanh nghiệp quản lý khách hàng, sản phẩm, đơn hàng và trách nhiệm xử lý thông qua quy trình tùy chỉnh, kết hợp phân hệ hội thoại và AI hỗ trợ tư vấn có kiểm soát, đặt trên nền tảng kiến trúc mở sẵn sàng cho định hướng hệ sinh thái vận hành tiếp thị -- bán hàng. Ba câu hỏi nghiên cứu trọng tâm gồm:
\begin{enumerate}
  \item Quy trình công việc tuần tự tùy chỉnh có phiên bản bất biến và hai chế độ bàn giao (chuyển trực tiếp và chờ xác nhận nhận) có biểu diễn được các cách tổ chức vận hành khác nhau, đồng thời xác định rõ bước xử lý và người chịu trách nhiệm tại từng thời điểm hay không?
  \item Cơ chế định tuyến kịch bản mẫu và AI soạn nháp dựa trên dữ liệu doanh nghiệp được cấp phép có giúp giảm thời gian và thao tác chuẩn bị phản hồi so với cách dùng lệnh trả lời nhanh (tính cả công kiểm tra, chỉnh sửa) mà vẫn bảo đảm chất lượng thông tin và không tự suy đoán dữ liệu thiếu hay không?
  \item Hệ thống có bảo đảm phân quyền theo vai trò, cô lập dữ liệu giữa các doanh nghiệp, kiểm soát cập nhật cạnh tranh và lưu vết lịch sử thay đổi đầy đủ trong các kịch bản thử nghiệm hay không?
\end{enumerate}

Bên cạnh ba câu hỏi nghiên cứu, khóa luận đặt một mục tiêu kỹ thuật về khả năng chịu tải: trên hai máy chủ 2 vCPU, gồm một máy 4 GB RAM chạy hạ tầng dữ liệu (PostgreSQL, Kafka, Redis) và một máy 16 GB RAM chạy các dịch vụ ứng dụng, hệ thống phải nâng thông lượng tiếp nhận từ mốc xuất phát khoảng 1 tin nhắn/giây lên 300 tin nhắn/giây duy trì, chịu đột biến 1.000 tin nhắn/giây, gửi 100 tin/giây ra kênh ngoài và phục vụ 500 yêu cầu/giây từ nhân viên mà không mất hoặc trùng tin, kể cả khi Kafka tạm ngừng (chi tiết tại Bảng~\ref{tab:load-targets}).

Về định hướng đánh giá, câu hỏi nghiên cứu thứ hai tập trung đo kiểm trên chế độ soạn nháp có người duyệt để bảo đảm đối chứng rõ ràng; chế độ tự gửi đòi hỏi bộ kiểm chứng an toàn và quy trình đối soát riêng biệt, không suy diễn từ kết quả thử nghiệm soạn nháp. Phạm vi thực nghiệm sử dụng môi trường triển khai cục bộ có thể tái lập cùng các kịch bản kiểm thử có kiểm soát và mô phỏng đối tác; dữ liệu mô phỏng được gắn nhãn tường minh và chưa được coi là giao dịch thực tế. Ranh giới phạm vi của khóa luận được phân định cụ thể trong Bảng~\ref{tab:scope-boundary}.

\begin{table}[htbp]
  \centering
  \small
  \setlength{\tabcolsep}{4pt}
  \renewcommand{\arraystretch}{1.2}
  \caption{Phân định phạm vi nghiên cứu của khóa luận}\label{tab:scope-boundary}
  \begin{tabularx}{\linewidth}{@{}
    >{\raggedright\arraybackslash}p{3.4cm}
    >{\raggedright\arraybackslash}X
    >{\raggedright\arraybackslash}p{3.8cm}@{}}
    \toprule
    Phân vùng phạm vi & Nội dung thuộc nhóm & Mức độ trong khóa luận \\
    \midrule
    Phạm vi cốt lõi của khóa luận &
    Quản lý đa doanh nghiệp (xét duyệt đăng ký và lời mời thành viên), khách hàng, sản phẩm, đơn nháp và xác nhận có điều kiện, quy trình tuần tự có phiên bản với hai chế độ bàn giao, điều phối hội thoại (tự động, chờ tiếp nhận, nhân viên xử lý), AI định tuyến và soạn nháp có kiểm chứng dữ liệu, kết nối kênh trực tiếp (Facebook Page, Telegram Bot) và bộ mô phỏng gắn nhãn. &
    Phân tích, thiết kế, cài đặt và kiểm chứng trong môi trường thực nghiệm có kiểm soát. \\
    Biên kiến trúc mở hướng hệ sinh thái &
    Chuẩn hóa định danh doanh nghiệp dùng chung, mô hình phân quyền hạt nhân, tách biệt biên dịch vụ kết nối kênh và luồng sự kiện miền bất đồng bộ. &
    Thiết kế sẵn cấu trúc nền tảng để mở đường kết nối các ứng dụng tiếp thị về sau. \\
    Ngoài phạm vi hiện tại &
    Các ứng dụng vệ tinh chuyên biệt cho sản xuất nội dung truyền thông và quản trị xuất bản trang mạng xã hội; nghiệp vụ thu tiền, hoàn tiền, kế toán đầy đủ, quản lý kho toàn diện, chấm công/lương, phân tích quảng cáo chuyên sâu, quy trình rẽ nhánh/song song và vận hành sản xuất quy mô lớn. &
    Thuộc định hướng phát triển mở rộng dài hạn; không cài đặt và nghiệm thu trong đồ án. \\
    \bottomrule
  \end{tabularx}
\end{table}

\section{Nội dung khóa luận}
Để trả lời ba câu hỏi nghiên cứu và đạt được mục tiêu đề ra, khóa luận triển khai bốn nhiệm vụ chính:
\begin{enumerate}
  \item \textbf{Phân tích bài toán và khảo sát giải pháp:} Hệ thống hóa các vấn đề vận hành từ tiếp thị đầu phễu đến tư vấn, đơn hàng và bàn giao nội bộ; đối chiếu các nền tảng hiện có và xác lập bộ yêu cầu chức năng cùng yêu cầu phi chức năng cho hệ thống.
  \item \textbf{Thiết kế kiến trúc phân tách trách nhiệm và biên mở rộng:} Xây dựng kiến trúc đa thành phần phân tách rõ vai trò giữa lớp giao diện, lớp nghiệp vụ trung tâm nắm quyền xác thực và cấp phép, lớp kết nối kênh giao tiếp ngoài và lớp thực thi AI chỉ trả về ứng viên, đồng thời giữ biên định danh và sự kiện mở cho định hướng hệ sinh thái tương lai.
  \item \textbf{Thiết kế mô hình dữ liệu và cơ chế đảm bảo nhất quán:} Thiết kế lược đồ dữ liệu đa doanh nghiệp, cơ chế kiểm soát cập nhật cạnh tranh theo phiên bản, hàng rào bảo vệ quyền điều khiển hội thoại, cơ chế lưu vết lịch sử và truyền thông sự kiện bền vững nhằm ngăn ngừa sai lệch dữ liệu và gửi lặp.
  \item \textbf{Triển khai và thực nghiệm đánh giá:} Cài đặt các phân hệ cốt lõi trên môi trường cục bộ có thể tái lập và tiến hành các kịch bản kiểm thử có kiểm soát về cô lập doanh nghiệp, bàn giao quy trình, thực thi AI bền vững và tiếp nhận sự kiện kênh để đánh giá kết quả theo ba câu hỏi nghiên cứu, đồng thời kiểm thử tải đối chiếu với mục tiêu chịu tải đã đặt ra.
\end{enumerate}

\section{Cấu trúc khóa luận}
Ngoài phần mở đầu, tài liệu tham khảo và phụ lục, nội dung chính của khóa luận được tổ chức thành năm chương:
\begin{itemize}
  \item \textbf{Chương~\ref{chap:introduction} -- Giới thiệu:} Trình bày bối cảnh thực tiễn, bài toán cần giải quyết, đối chiếu các nhóm giải pháp hiện có, định hướng khác biệt trên kiến trúc mở, mục tiêu, phạm vi và bố cục khóa luận.
  \item \textbf{Chương~\ref{chap:background} -- Cơ sở lý thuyết:} Trình bày các khái niệm nền tảng về dữ liệu đa doanh nghiệp, tách biệt nháp và xác nhận, quy trình bất biến, phân tích chuyên sâu các giải pháp liên quan cùng các công nghệ và mẫu kiến trúc áp dụng.
  \item \textbf{Chương~\ref{chap:design} -- Phân tích và thiết kế hệ thống:} Mô tả các tác nhân, đặc tả yêu cầu chức năng và phi chức năng, kiến trúc tổng thể, thiết kế cơ sở dữ liệu và chi tiết các luồng xử lý nghiệp vụ cốt lõi.
  \item \textbf{Chương~\ref{chap:evaluation} -- Triển khai và thực nghiệm:} Trình bày môi trường cài đặt, thiết kế kịch bản thực nghiệm, chi tiết cài đặt các chức năng chính cùng kết quả đánh giá và giới hạn kết luận.
  \item \textbf{Chương~\ref{chap:conclusion} -- Kết luận:} Tổng kết các kết quả đạt được đối chiếu với ba câu hỏi nghiên cứu, chỉ ra hạn chế hiện tại và đề xuất hướng phát triển mở rộng hệ sinh thái.
\end{itemize}
Bên cạnh năm chương chính, \textbf{Phụ lục~\ref{app:database}} trình bày chi tiết cấu trúc các bảng cơ sở dữ liệu cốt lõi và tiến trình phát triển lược đồ của hệ thống.
